\section{Introduction}
\label{sec:introduction}

Security operations and digital-forensic investigations reconstruct multi-step attacks from
temporally dispersed endpoint and network observations~\cite{alertfatigue2025,nist80086,breitinger2025sok}.
Isolated EDR or NDR records may appear benign alone yet become material when linked by time,
entity, session, and process relations. Reconstruction is not only a clustering problem:
investigators need a contract that fixes who owns episode membership, how late evidence is
admitted, how collector health differs from ``no detection,'' and how every revision leaves an
independently verifiable terminal outcome~\cite{casey2011digitalevidence,carrier2005filesystem,vanini2024timeanchors}.

Prior alert-correlation and provenance systems motivate multi-stage context using rules, graphs,
and learned representations~\cite{ning2002attackscenarios,holmes2019,unicorn2020,nodlink2024,kairos2024,captain2025}.
Those lines do not remove the need for fixed ownership of membership, immutable revision history,
and replayable audit semantics under incomplete observability. This submission evaluates
\trustepisode as a \emph{deterministic, auditable reconstruction contract} for EDR/NDR evidence.
A compact scorer and calibrator remain optional schema interfaces and are not empirically
validated here.

\paragraph{Research questions.}
\begin{enumerate}[leftmargin=*,itemsep=1pt,topsep=2pt]
  \item[\textbf{RQ1}] Under identical sealed EDR/NDR evidence, does TrustEpisode reproduce
        byte-identical membership, revision history, late-evidence records, and digests on all
        60 malicious runs?
  \item[\textbf{RQ2}] Do open, finalized, late-successor, closed, beyond-horizon, duplicate, and
        all seven AssemblyFailureRecord codes yield explicit immutable audit outcomes without
        rewriting earlier history?
  \item[\textbf{RQ3}] When a required action is absent from a reconstructed episode, can the
        pipeline localize loss across the closed A--G taxonomy, and is that taxonomy instantiated
        on sealed live evidence?
  \item[\textbf{RQ4}] Which relation, hub, merge, and late-revision mechanisms does the live
        M1--M6 cohort exercise, and which require targeted conformance cards?
  \item[\textbf{RQ5}] What portion of observed formation gain comes from replacing fixed bins
        with sliding sessionization, and does relation-aware logic provide additional measurable
        benefit on the sealed cohort?
\end{enumerate}

\paragraph{Contributions.}
\begin{enumerate}[leftmargin=*,itemsep=1pt,topsep=2pt]
  \item a single-owner immutable reconstruction contract: Episode Synthesis alone owns
        membership; late evidence yields successors or LateEvidenceRecords; deadlines do not
        extend; terminal history is append-only;
  \item replayable forensic audit outcomes with exactly one terminal AssemblyOutcome per
        revision, canonical references, SHA-256 digests, and a complete seven-code failure matrix;
  \item a stage-localized A--G evidence-loss taxonomy, conformance-tested across classes and
        instantiated on sealed M2 as \texttt{D\_token\_mapping\_failure};
  \item a fair negative-result evaluation: legacy fixed bins underperform sliding sessionization,
        while fair sliding EB1--EB4 show no observed difference on 60 sealed runs; advanced
        guards require targeted cards not mixed into primary effectiveness;
  \item a determinism boundary on all 60 sealed malicious runs: canonical ordered replay and
        duplicate retries are stable; raw delivery-order permutations are not order-robust
        without re-canonicalization; formation is single-writer per run;
  \item a contract-ablation comparison (AB0--AB3) isolating which forensic failure modes mutable,
        noncanonical, and silent-failure stores admit relative to TrustEpisode.
\end{enumerate}
